\documentclass[11pt]{article}

\usepackage[margin=1in]{geometry}
\usepackage{amsmath,amssymb}
\usepackage{enumitem}
\usepackage{xcolor}
\usepackage[colorlinks=true,linkcolor=blue,urlcolor=blue,citecolor=blue]{hyperref}
\usepackage{microtype}
\usepackage{parskip}
\setlength{\emergencystretch}{2em}

\newcommand{\msnumber}{JME-D-26-00443}
\newcommand{\commenttext}[1]{\par\noindent\textbf{Comment.} \textit{#1}\par}
\newcommand{\response}[1]{\par\noindent\textbf{Response.} #1\par}

\begin{document}

\begin{center}
{\Large\bfseries Response to the Editor and Referees}\par
\vspace{0.5em}
{\large Asset Value and Securitization under Heterogeneous Beliefs}\par
\vspace{0.25em}
Manuscript \msnumber\par
\vspace{0.25em}
Oliver Pardo
\end{center}

\bigskip

I thank the editor and the two referees for identifying the failure of the original uniqueness argument. The reports revealed that this failure is not merely a gap in the proof. It reflects an economically meaningful feature of the model: the no-tranching benchmark has a unique equilibrium, whereas tranching can create multiple equilibrium prices. The revised paper establishes this result formally, presents a two-tranche debt--equity example with three equilibria, and develops the equilibrium selection needed for the comparative statics. It then re-examines the downstream results and distinguishes those that depend on the selection from those that hold for every equilibrium.

The main changes are as follows.

\begin{enumerate}[leftmargin=2em]
\item Proposition 1 now establishes existence of least and greatest waterfall-equilibrium prices, not uniqueness. A new corollary shows that the no-tranching benchmark has a unique equilibrium, whereas finite tranching can create multiplicity. The revised paper adapts the referee's exact three-fixed-point example by merging its two most senior tranches; the resulting two-tranche debt--equity structure preserves all three fixed points.
\item The comparative statics use the least fixed point, obtained by iteration from zero continuation values. A new subsection on equilibrium selection and Proposition 2 show that this selection is the limit of backward-induction prices in finite-horizon economies with zero terminal resale value.
\item The price results are separated from the Eyster--Piccione partition structure. Pricing permits any finite collection of full-support subjective Markov matrices; statistical consistency and partition refinement are imposed only where the realized-payoff results require them.
\item The results on refinement, issuer optimality, individual valuations, fundamentals, and repeated securitization have been restated and reproved under the new equilibrium treatment. Results that apply to every waterfall equilibrium are identified explicitly.
\item When at least one dividend is positive, Proposition 8 gives a necessary and sufficient condition for the state-contingent price equation to have a nonnegative solution: all eigenvalues of the relevant matrix must lie strictly inside the unit circle. When this condition holds, Proposition 9 shows that the resulting price vector bounds every waterfall-equilibrium price from above.
\item The realized-payoff section now fixes an exogenous strict total priority ordering over theories and proves a new aggregate result. Under statistical consistency, the aggregate realized excess payoff from unit tranche purchases is nonpositive. Under the equilibrium selection used in the paper, this aggregate payoff weakly decreases under refinement and decreases strictly whenever the selected asset price rises.
\item The introduction now engages directly with the securitization literature identified by the referee and states both the contribution of the long-lived-asset feedback and the restrictions of the model. The revised paper also distinguishes realized-payoff comparisons from welfare rankings.
\end{enumerate}

References below are to the revised manuscript.

\section*{Response to the Editor}

\commenttext{The failure of uniqueness in Proposition 1 must be addressed at the theorem level. Any selection must be defined precisely, justified economically, and carried through all price and issuer-optimality results.}

\response{I agree. The old Blackwell argument and the uniqueness claim have been removed. Proposition 1 now proves that the fixed-point set is nonempty and has least and greatest elements. The proof uses monotonicity and continuity, constructs an explicit upper bound on prices that is preserved by the price operator, and iterates the operator from zero and from that upper bound. Immediately afterward, the paper presents a simplified version of the referee's exact economy. Merging the two most senior tranches produces a two-tranche debt--equity structure while preserving all three fixed points. A new corollary proves that the no-tranching operator remains a contraction and therefore has a unique fixed point. The counterexample then establishes that tranching itself can create multiplicity.

Throughout the remaining analysis, $q_{\mathcal T}$ is defined as the least fixed point. A new subsection on equilibrium selection contains Proposition 2, which shows that this price is the limit of the unique backward-induction prices in finite-horizon versions of the model. Lemma 1 compares the price sequences generated by a tranching and its refinement when both price operators are iterated from zero continuation values. Its conclusion is therefore explicitly selection-dependent. Proposition 3 and the definition of issuer optimality are likewise stated under the least-equilibrium selection. By contrast, Proposition 4, which compares any equilibrium price with individual one-step valuations, and Proposition 5, which bounds any equilibrium price below by perceived fundamentals, are stated for every element of the fixed-point set.

Proposition 2 gives the selection an economic interpretation. Starting from zero terminal continuation value and solving backward in the finite-horizon economy with $H$ remaining dividend dates produces the $H$-th iterate of the price operator. These finite-horizon prices increase to the least stationary fixed point. Thus, the selection is not imposed only as an order-theoretic convention; it is the stationary limit of a sequence of uniquely determined finite-horizon prices.}

\commenttext{Correct the asserted existence of a nonnegative state-contingent price and state the comparison with waterfall tranching under sufficient conditions.}

\response{Done. The revised text defines the nonnegative matrix
\[
A(x,y)=\frac{\max_{\mathcal F}Q_{\mathcal F}(x,y)}{R(x)}.
\]
Proposition 8 shows that, when at least one dividend is positive, the state-contingent price equation $\bar q=b+A\bar q$ has a nonnegative solution if and only if all eigenvalues of $A$ lie strictly inside the unit circle. Necessity follows by multiplying this equation by a positive left eigenvector associated with the dominant eigenvalue of $A$. When the condition holds, the solution is unique and equals $\bar q=(I-A)^{-1}b=\sum_{n=0}^{\infty}A^n b$. Proposition 9 then proves that every waterfall-equilibrium price $q$ generated by a finite tranching satisfies $q\leq\bar q$.}

\commenttext{Specify an allocation or tie-breaking rule and make the return definitions, tables, and examples consistent with it.}

\response{Done. The revised realized-payoff section keeps the set of maximizing theories $\mathbf B(\tau,x)$ for pricing and fixes an exogenous strict total priority ordering over theories, independent of the current state and the tranche. The selected buyer $\beta_{\mathcal T}(\tau,x)$ is the highest-priority element of $\mathbf B(\tau,x)$. The definition of realized excess payoffs, the payoff calculations, and the table identifying the selected buyers all use this rule. Prices do not depend on the priority ordering, while realized excess payoffs can depend on it.

The state-independent priority ordering also closes a logical gap in the individual-payoff comparison. When one theory refines every other theory, the set of maximizing theories is constant on each block of the finest theory. The selected buyer is therefore constant on that block as well. The revised paper now gives the blockwise calculation explicitly and proves, tranche by tranche, that the finest theory has zero realized excess payoff and every coarser theory has a nonpositive one.}

\commenttext{Position the contribution relative to the literature identified by Reviewer 2 and explain what is gained from the infinite-horizon Harrison--Kreps structure and the long-lived underlying asset.}

\response{The introduction and literature review now discuss Guo, Wang, and Wei (2014), Broer (2018), Ellis, Piccione, and Zhang (2022), and Iachan, Nenov, and Simsek (2021). The revised comparison does not claim that the basic heterogeneous-beliefs tranching mechanism is new. It identifies the distinct object studied here: securitization changes the endogenous resale value of a long-lived asset; that resale value determines future tranche payoffs; and the resulting loop can generate multiple stationary prices. This feedback, together with the evaluation of long-run realized excess payoffs under an objective transition process, is now presented as the paper's contribution. The literature review and conclusion also make the paper's limitations explicit: the model assumes risk-neutral agents, takes the waterfall class and the collection of beliefs as given, and does not endogenize risk pooling.}

\commenttext{Separate pricing results that hold for arbitrary Markov beliefs from return results that require statistical consistency or partition refinement, and clarify the persistence of conditionally incorrect beliefs.}

\response{Done. The beliefs subsection now begins with an arbitrary finite collection of full-support subjective Markov matrices. It then introduces the Eyster--Piccione partition structure as an additional assumption used primarily for realized excess payoffs. The text states explicitly that statistical consistency is unconditional calibration, not conditional correctness, and that the theories are maintained reduced-form forecasting rules rather than Bayesian posteriors in a correctly specified common model. Endogenous learning or switching among theories is outside the equilibrium analyzed in the paper.}

\commenttext{Clarify the economic interpretation and limitations of the securitization structure and, if feasible, provide a finite-horizon robustness analysis or illustration.}

\response{A new subsection on equilibrium selection provides the finite-horizon illustration without replacing the stationary model. Proposition 2 shows that, in an economy with $H$ remaining dividend dates, zero terminal continuation value, and the same one-period waterfall menu at every date, backward induction is exactly iteration of the price operator from zero. The proof of Lemma 1 also shows that refinement weakly raises the price in every such finite-horizon economy. These prices converge to the selected stationary equilibrium as $H$ grows. The finite-horizon exercise preserves the paper's one-period waterfall claims on next-period continuation values. It therefore does not establish that the results extend to ABS backed only by dividends paid over a fixed number of periods, to multi-period tranches, or to alternative default structures.}

\commenttext{Address the remaining terminology, CDF, finiteness, iterated-securitization, and expositional comments.}

\response{All have been addressed. The introduction now refers to a finite set of belief types rather than a finite number of investors. The CDF inequalities in the motivating example are restricted to the intervals on which they are strict. ``Non-arbitrage'' has been replaced by ``no-arbitrage.'' The iterated-securitization numerical term and budget-balance identity have been corrected. The appendix now proves that an iterated securitization structure and its flattened one-shot waterfall tranching generate the same complete set of equilibrium prices for the underlying asset; consequently, their least equilibrium prices coincide.}

\section*{Response to Reviewer 1}

\subsection*{Major comment 1: Proposition 1 and the failed discounting inequality}

\commenttext{The aggregate value of separately priced tranches can respond to a constant shift by more than that shift. The proof of uniqueness therefore has a gap, and later results that rely on uniqueness must be reconsidered.}

\response{I agree. The discounting claim was false and has been deleted. Proposition 1 now establishes existence of least and greatest fixed points by monotone iteration on a compact invariant order interval. The paper no longer claims uniqueness of a waterfall equilibrium. It proves separately that the no-tranching benchmark has a unique equilibrium and that tranching can create multiplicity. A new subsection defines the least equilibrium as the selection used for the refinement and issuer-optimality results, and Proposition 2 shows that this selection is the limit of the unique backward-induction prices in finite-horizon versions of the model. The paper also identifies separately the propositions that hold for every equilibrium.}

\subsection*{Major comment 2: Why impose partition-generated beliefs?}

\commenttext{The main pricing mechanism appears to require only finitely many subjective Markov transition matrices, whereas statistical consistency and partition refinement matter for particular return results.}

\response{The revised paper adopts exactly this separation. The pricing sections permit an arbitrary finite collection of full-support stochastic matrices. Full support is used to propagate a strict price increase from one state to every state. Neither the objective matrix $P$, its stationary distribution $\mu$, nor partition-generated beliefs are otherwise needed for the price operator.

The realized-payoff section then imposes the Eyster--Piccione structure explicitly. The common stationary distribution is enough for Proposition 6, which establishes the sign of the aggregate per-unit excess payoff. Under the least-equilibrium selection, a new corollary combines Proposition 6 with the price-comparison lemma to show that this aggregate payoff weakly decreases under refinement. The stronger conclusion that a finest theory earns zero and every coarser theory earns a nonpositive realized excess payoff uses the partition order, the fixed priority ordering, and the blockwise argument underlying Proposition 5 in Eyster and Piccione. The revised manuscript provides that argument explicitly.}

\subsection*{Minor comment 3: Finite investors versus finite theories}

\response{The introduction now says that risk-neutral investors have a finite set of belief types. It no longer attributes a substantive role to a finite number of investors.}

\subsection*{Minor comment 4: Strict CDF inequalities}

\response{Corrected. At state $m$, the revised text states $G_{\mathcal B}(t\mid m)<G_{\mathcal A}(t\mid m)$ only for $q(l)\leq t<q(m)$ and $G_{\mathcal A}(t\mid m)<G_{\mathcal B}(t\mid m)$ only for $q(m)\leq t<q(h)$. It also states that the CDFs coincide below $q(l)$ and at or above $q(h)$.}

\section*{Response to Reviewer 2}

\subsection*{Major comment 1: Exact multiplicity and equilibrium selection}

\commenttext{The waterfall equilibrium need not be unique. The report provides three exact fixed points for one tranching and notes that one is also the no-tranching price.}

\response{I agree and thank the referee for the counterexample. The revised manuscript reproduces the transition matrix, interest rate, dividend vector, and all three exact fixed points. It merges the two most senior tranches in the referee's construction into senior debt with face value 17. The resulting two-tranche debt--equity structure preserves all three fixed points and shows that the three equilibria correspond to different sets of states in which residual equity pays. Proposition 1 now establishes extremal fixed points rather than uniqueness. A new corollary proves that the no-tranching benchmark has a unique equilibrium, so the example establishes that tranching can create multiplicity rather than merely revealing multiplicity already present in the benchmark. All later comparative statics have been audited for their dependence on the least-equilibrium selection.}

\subsection*{Major comment 2: Placement}

\commenttext{The contribution should be compared more carefully with Guo, Wang, and Wei; Broer; Ellis, Piccione, and Zhang; and Iachan, Nenov, and Simsek. The long-lived asset and feedback into resale prices appear to be the natural distinction.}

\response{The revised literature review now engages directly with all four references. The introduction makes the comparison in the terms suggested by the referee. Prior work already establishes that diverse beliefs can make tranching valuable because different buyers price different claims. The paper's distinct object is the stationary feedback between securitization, the resale value of the long-lived underlying asset, and future tranche payoffs. The revision also states that this feedback is responsible for the newly acknowledged multiplicity and supports the long-run realized-payoff analysis.}

\subsection*{Major comment 3: Returns and welfare}

\commenttext{The realized-return analysis is potentially more novel but underdeveloped. A sharper characterization would help; any welfare discussion should engage with Brunnermeier, Simsek, and Xiong.}

\response{The realized-payoff section now contains Proposition 6, which goes beyond the examples. Under statistical consistency and any fixed priority ordering, the aggregate realized excess payoff from unit tranche purchases is
\[
\Delta_{\mathrm{agg}}^{\mathcal T}
=\sum_x\mu(x)\left[d(x)-(R(x)-1)q_{\mathcal T}(x)\right]\leq0.
\]
A new corollary combines this identity with the price comparison under refinement: if $\mathcal S$ refines $\mathcal T$, then $\Delta_{\mathrm{agg}}^{\mathcal S}\leq\Delta_{\mathrm{agg}}^{\mathcal T}$, with strict inequality whenever the selected price rises. The proposition concerns the sum of excess payoffs from one-unit tranche purchases, rather than a portfolio-investment-weighted rate of return. Individual effects remain ambiguous: the existing examples show both that securitization can lower one type's realized excess payoff without raising another's and that it can make another type's negative excess payoff less negative by shifting a worse claim.

The revised paper does not infer welfare from these payoff comparisons. The conclusion explains that risk neutrality rules out risk-sharing gains in the present model and that evaluating whether such gains under risk aversion would outweigh the potential welfare costs of belief-driven speculation requires a richer model. It identifies the belief-neutral efficiency criterion of Brunnermeier, Simsek, and Xiong (2014) as one possible approach.}

\subsection*{Major comment 4: Maturities and the securitization structure}

\commenttext{The one-period claims are written on the next-period value of an infinitely lived asset. A simple finite-horizon illustration would clarify which results rely on the stationary resale structure.}

\response{The new subsection on equilibrium selection supplies this illustration. Proposition 2 shows that, with $H$ remaining dividend dates, zero terminal continuation value, and the same one-period waterfall menu at each date, backward induction gives the $H$-th price vector generated by iterating the waterfall operator from zero continuation values. The proof of Lemma 1 shows that refinement raises prices weakly in every such finite-horizon economy, and Proposition 2 shows that these prices converge to the selected infinite-horizon price. The revision therefore separates the finite-horizon comparison from the stationary fixed-point feedback that creates multiplicity and repeated amplification. The construction does not cover alternative maturity or default structures, and claims only on a currently known dividend cannot generate the mechanism because their payoff has no future-state variation.}

\subsection*{Minor comment 1: Existence of a nonnegative state-contingent price}

\response{Corrected as described in the response to the editor. When at least one dividend is positive, Proposition 8 gives a necessary and sufficient condition for the state-contingent price equation $\bar q=b+A\bar q$ to have a nonnegative solution. Proposition 9 proves the upper-bound comparison for every waterfall equilibrium whenever this solution exists.}

\subsection*{Minor comment 2: Allocation between tied buyers}

\response{A fixed exogenous strict total priority ordering over theories now determines $\beta_{\mathcal T}$ in the definition of realized excess payoffs and in every relevant table and example. The ordering is independent of the current state and the tranche. This both makes the allocation explicit and ensures that, when a finest theory exists, selected-buyer sets are unions of blocks of that theory.}

\subsection*{Minor comment 3: Beliefs and learning}

\response{The pricing analysis now permits arbitrary finite full-support Markov beliefs. Partition-generated statistical consistency is introduced only for the realized-payoff analysis. The paper also clarifies that the maintained theories are persistent reduced-form forecasting rules and that endogenous learning is not part of the model; the learning-based crash mentioned in the conclusion is explicitly labeled as a possible extension.}

\subsection*{Minor comment 4: Terminology}

\response{``Non-arbitrage condition'' has been replaced by ``no-arbitrage condition.''}

\subsection*{Minor comment 5: Iterated-securitization example}

\response{Corrected. In the iterated-securitization example, the final contribution to the asset price at state $d$ is now $\tfrac12[q(a)-q(b)]$. During the same audit, I also corrected the budget-balance identity that sums the payoffs of a node's successor securities and strengthened the equivalence proposition to show that an iterated structure and its flattened one-shot waterfall tranching generate the same complete set of equilibrium prices for the underlying asset, rather than implicitly relying on uniqueness.}

\bigskip

I hope the revised manuscript resolves the mathematical, interpretive, and positioning issues raised in the reports. I am grateful for the counterexamples and for the suggestions that led to a more accurate and substantially stronger paper.

\end{document}
